\chapter{Algorithm and pseudocode}
\label{app:algorithm}

\section{Validation algorithm}

\begin{lstlisting}[language=Julia]
function validate_project(config)
    project = load_project(config)
    data = load_network(project)

    check_schema(data)
    check_mode_and_congestion_keys(project, data)
    check_physical_policy_pairs(project, data)
    check_flow_and_exposure_accounting(data)

    basis = reconstruct_route_basis(project, data;
                                    include_fixed=false)
    require(basis.spectral_radius < 1)
    require(basis.margin_error <= tolerance)
    require(basis.edge_error <= tolerance)

    closure = build_full_welfare_closure(project, data, basis)
    require(condition_number(closure) <= condition_limit)
    return validation_summary(project, data, basis, closure)
end
\end{lstlisting}

Validation builds only the objects needed for the welfare closure. It does not
construct the full fixed-route incidence matrix.

\section{Economic-geography welfare algorithm}

\begin{lstlisting}[language=Julia]
function all_welfare_effects(project)
    data = load_network(project)
    basis = reconstruct_route_basis(project, data;
                                    include_fixed=false)

    coefficients = spatial_coefficients(project.parameters)
    J0 = no_congestion_jacobian(data, coefficients)
    q = welfare_gradient(data.income_share, coefficients)
    B = aggregate_cost_loading(data, coefficients)

    transport_F = transport_closure(project, data, basis,
                                    flexible_routes=true,
                                    flexible_modes=true,
                                    edge_congestion=true,
                                    terminal_congestion=true)
    JF = J0 + transport_F.feedback

    b_primitive = primitive_policy_forcing(B, transport_F, basis)
    adjoint = solve(transpose(JF), q)
    elasticities = transpose(adjoint) * b_primitive

    return directed_and_physical_results(elasticities, data, project)
end
\end{lstlisting}

The implementation uses sparse matrices where possible. The schematic
\code{solve} represents a factorization and backsolve, not an explicit inverse.

\section{Full decomposition algorithm}

\begin{lstlisting}[language=Julia]
function decompose(project)
    data, basis = build_route_basis(project; include_fixed=true)
    J0, q, B = build_spatial_system(project, data)

    NC = closure(J0, no_congestion,
                 flexible_routes, flexible_modes)
    NT = closure(J0, edge_congestion,
                 flexible_routes, flexible_modes)
    F  = closure(J0, full_congestion,
                 flexible_routes, flexible_modes)
    FM = closure(J0, full_congestion,
                 flexible_routes, fixed_modes)
    FR = closure(J0, full_congestion,
                 fixed_routes, flexible_modes)

    b_realized = realized_policy_forcing(B, basis)
    b_primitive = primitive_policy_forcing(B, F, basis)

    E = adjoint_gains((NC, NT, F, FM, FR), q, b_realized)
    Eprimitive = adjoint_gain(F, q, b_primitive)

    check_inverse_gap_identities(E, (NC, NT, F, FM, FR))
    channels = analytical_mixed_updates(F, FM, FR)
    check_channel_reconstruction(channels)
    check_hulten_accounting(E, Eprimitive, traffic)

    return write_decomposition(E, Eprimitive, channels)
end
\end{lstlisting}

\section{Route reconstruction}

For distinct source and destination margins:
\begin{lstlisting}[language=Julia]
K = continuation_matrix(edge_traffic, exposure)
T = inverse(I - K)
Xod = Diagonal(source_margin) * T * Diagonal(absorption)

require(maximum(abs.(sum(Xod, dims=2) - source_margin)) <= tol)
require(maximum(abs.(sum(Xod, dims=1)' - destination_margin)) <= tol)
require(edge_incidence(Xod, K) matches observed_edge_traffic)
\end{lstlisting}
The code uses linear solves rather than materializing \code{inverse(I-K)} when
possible. The notation highlights the resolvent identity.

\section{Primitive forcing}

Let \(p\) index active edge-mode pairs and let \(D\) select policy pairs.
For the full closure:
\[
 X_{\theta}
 =L(C^{route}S_{\mathrm{agg}}D)
  +p_m(D-LS_{\mathrm{agg}}D),
\]
where \(p_m=-\eta\) for the choice logsum. Congestion quantities respond as
\[
 q_\theta
 =(I-AX_q)^{-1}AX_\theta.
\]
The aggregate edge-cost change is
\[
 h_\theta=S_{\mathrm{agg}}(D+Gq_\theta),
\]
and the spatial forcing is
\[
 b^\theta=B_\kappa h_\theta.
\]
These operations are implemented in \code{primitive\_forcing}. Shared terminal
congestion can therefore make primitive pass-through nonlocal.

\section{Complexity}

Let \(N\) be nodes, \(E\) active directed edges, \(P\) active edge-mode pairs,
and \(Q\) congestion states. The welfare path requires route reconstruction,
a \(2N\)-scale spatial factorization, a \(Q\)-scale transport solve, and sparse
products for all policy pairs. The full route decomposition additionally
stores origin-destination edge incidence, which can scale as \(N^2E\).

The validator reports a preflight memory estimate. For large assignment
networks, use \code{analyze} unless the fixed-route comparison is essential and
the incidence object is computationally feasible.
